\documentclass[../main.tex]{subfiles}

\begin{document}

Lean es un asistente de demostración y lenguaje de programación funcional.
Está basado en la teoría de tipos dependientes, gracias a esto y al 
isomorfismo de Howard-Curry, podemos usar el mismo lenguaje para escribir
programas y demostraciones. Existe una librería matématica Mathlib para Lean4, que
contiene una gran cantidad de demostraciones clásicas,
estas van desde teoría de conjuntos hasta topología.

Sus creadores también hacen énfasis en su uso como lenguaje de programación
funcional, ya que es un lenguaje de propósito general, aunque nosotros
nos centraremos en su uso como asistente de demostración.

\section{comparación de Lean con otros asistentes de demostración}

Lean4 y Coq son dos de los asistentes de demostración más populares.
Son parecidos en que ambos usan teoría de tipos, en particular una
versión de la teoría de tipos dependientes llamada CIC 
(Calculus of Inductive Constructions). En ambos sistemas las proposiciones
son tipos que viven en un universo especial, al margen de la jerarquía de
tipos. Demostrar una proposición es construir un término del tipo
correspondiente a la proposición. Esto es, en ambos, la proposición
\(\forall x \in \mathbb{N}, x + 0 = x\) es un tipo, y la veracidad,
o no, de esta dependerá de si este tipo está habitado, por lo que demostrar
una proposición es construir un término de este tipo. Podemos ver esto
como una versión del isomorfismo de Howard-Curry.

\section{Lean 3 vs Lean 4}

\end{document}